\documentclass[11pt,a4paper]{ltjsarticle}
% Common preamble; compile the four documents with LuaLaTeX.
\usepackage{fix-cm}
\usepackage[margin=25mm]{geometry}
\usepackage{luatexja-fontspec}
\setmainjfont{HaranoAjiMincho-Regular.otf}[BoldFont=HaranoAjiMincho-Bold.otf]
\setsansjfont{HaranoAjiGothic-Medium.otf}[BoldFont=HaranoAjiGothic-Bold.otf]
\usepackage{amsmath,amssymb,amsthm,mathtools,booktabs,enumitem}
\usepackage[unicode,hidelinks]{hyperref}
\setlength{\emergencystretch}{2em}
\setlist{nosep}
\newtheorem{theorem}{定理}[section]
\newtheorem{proposition}[theorem]{命題}
\newtheorem{lemma}[theorem]{補題}
\newtheorem{corollary}[theorem]{系}
\theoremstyle{definition}
\newtheorem{definition}[theorem]{定義}
\newtheorem{remark}[theorem]{注意}
\newtheorem{example}[theorem]{例}
\newtheorem{problem}[theorem]{問題}
\newcommand{\C}{\mathbb C}
\newcommand{\Z}{\mathbb Z}
\newcommand{\Q}{\mathbb Q}
\newcommand{\Cat}{\mathcal A}
\newcommand{\DD}{\mathcal D}
\newcommand{\OO}{\mathcal O}
\newcommand{\TT}{\mathcal T}
\newcommand{\PP}{\mathcal P}
\newcommand{\SSer}{\mathbb S}
\newcommand{\Db}{D^{\mathrm b}}
\newcommand{\Hom}{\operatorname{Hom}}
\newcommand{\End}{\operatorname{End}}
\newcommand{\Ext}{\operatorname{Ext}}
\newcommand{\RHom}{\operatorname{RHom}}
\newcommand{\qgr}{\operatorname{qgr}}
\newcommand{\Gr}{\operatorname{Gr}}
\newcommand{\QGr}{\operatorname{QGr}}
\newcommand{\Tor}{\operatorname{Tor}}
\newcommand{\coh}{\operatorname{coh}}
\newcommand{\thick}{\operatorname{thick}}
\newcommand{\gldim}{\operatorname{gldim}}
\newcommand{\GKdim}{\operatorname{GKdim}}
\newcommand{\rk}{\operatorname{rk}}
\newcommand{\id}{\operatorname{id}}
\newcommand{\Hilb}{\operatorname{Hilb}}
\newcommand{\versiondate}{2026年11月17日・v14}

\renewcommand{\versiondate}{2026年11月17日・v14}
\newcommand{\Nnum}{N(\mathcal A)_{\Q}}
\newcommand{\Runi}{\mathcal R_{\mathrm u}}
\hypersetup{pdftitle={正規化された圏論的Fano指数}}
\title{正規化された圏論的Fano指数\\自己同型の混入・数値的標準類・曲面認識}
\author{}
\date{\versiondate\\研究原稿：一般の非可換Kobayashi--Ochiai型認識は未解決}
\begin{document}
\maketitle
\begin{abstract}
反Serre関手の根に曲面の自己同型を混ぜると，線束としての反標準類の整除性ではなくnorm条件が現れる．
従って根の存在だけを古典的Fano指数と同一視できない．この問題を具体例と一般のnorm公式で説明し，
数値Grothendieck群全体のunipotenceと，構造対象の軌道上だけのunipotenceを区別する．
滑らかな複素Fano多様体では両正規化で可能な次数は通常の指数の約数となるが，許される自己同値は異なる．
標識変更による正規化の障害はPicard群の巡回群コホモロジーに属する．
既存の曲面型擬格子の枠組みでは，unipotentな反Serre根の存在は標準類の整除性と同値になる．
defectが非正で正の標準次数を持つ幾何的擬格子では根の次数は3以下であり，
次数3と2の場合の格子はそれぞれ射影平面と二次曲面の既知の数値型に帰着する．
これは圏の認識ではない．別に，軌道Euler多項式から$r^2\mid d$，偶数$r$では$2r^2\mid d$を導く．
次数9の整三次normal境界では，正規化と既存の漸近消滅を用いて全軌道$\Ext^1$消滅と飽和性を得る．
さらに，正規化しない数値根を線束移動と有限位数のNS等長変換の半直積で分類し，\(r\mid K^2\)と非正の周期的Euler欠損を示す．
滑らかなdel Pezzo曲面の非正規化exact指数の次数別最大値を求め，次数2から6では反標準次数を達成する実在の例を与える．
また三次元ではSerre対称性とNewton展開だけで正規化根の整数値条件を得る．
先行研究の偏極・曲面型格子・root pair・geometric helixの定義を比較し，
既知分類の言い換え，本文で直接証明する帰結，未解決の圏論的持上げとfull性を区別する．
\end{abstract}
{\small\tableofcontents}

\section{問題の位置付けと二種類の正規化}
\subsection{自己同型を混ぜた根は，古典的指数ではない}\label{subsec:warning}
滑らかなFano多様体$X$の指数は$\omega_X^{-1}=H^{\otimes r}$となる最大の正整数$r$である．
しかし$\coh X$の自己同値には線束によるtensorだけでなく$X$の自己同型も入る．
Bondal--Orlov\cite[Theorem 3.1]{BO}に基づき$s=L\otimes f^*(-)$と書くと，
\begin{equation}\label{eq:intro-norm}
 s^r\simeq(-)\otimes\omega_X^{-1}
 \quad\Longleftrightarrow\quad
 f^r=\id_X,\quad\bigotimes_{j=0}^{r-1}(f^j)^*L\simeq\omega_X^{-1}.
\end{equation}
右辺はnorm条件であって，$L^{\otimes r}\simeq\omega_X^{-1}$ではない．
例えば次数6 del Pezzoの三点巡回と$L=\mathcal O(H-E_1)$は三乗根を，
$\mathbb P^1\times\mathbb P^1$の因子交換と$L=\mathcal O(1,0)$は四乗根を与える
（例\ref{ex:non-normalized}）．いずれもheart上exactかつAZ ampleなので，
exact性・ample性を付け加えるだけでもこのずれは除けない．
\emph{根の存在だけを用いた概念が無意味なのではなく，それは通常のFano指数とは別の，より広い偏極概念である．}
本稿の正規化は古典的指数を回復するための定義案であり，既存の定義の誤りを訂正するものではない．

\subsection{先行研究と本稿の課題}\label{subsec:literature}
Kuznetsov--Smirnov\cite[Definition 2.1, Example 2.2]{KS}は飽和三角圏で
$s^r\simeq\SSer^{-1}[a]$を指数$r$の偏極と定める．heart，AZ ample性，数値的unipotenceは要求しない．
本稿はこの根を正規化する問題を扱い，根そのものの定義を新設したとは主張しない．
古典的な最大指数の比較にはKachi--Koll\'ar\cite[Corollary 2(1)]{KK}を用いる．

Artin--Van den Berghのtwisted homogeneous coordinate ringの枠組みを拡張した
Keeler\cite[Theorems 1.2, 1.3]{Keeler}では，自己同型の数値的\emph{quasi-unipotence}と
有限個のtwistの和のample性が$\sigma$-amplenessを特徴付ける．
有限位数の非自明な固有値はquasi-unipotentなので，その理論は上の例を排除しない．
本稿のunipotenceによる正規化とは条件も目的も異なる．

数値的側面では，de Thanhoffer de Volcsey--Van den Bergh\cite[\S3.2--3.3]{dTVdB}の
surface-type Serre latticeと，Kuznetsov\cite[\S3, Definition 5.3]{Kuz}のsurface-like pseudolattice，
標準類，defectを基礎にする．これらにおける\emph{Serre作用}のunipotenceを，
その\emph{根}のunipotenceと混同しない．
本稿の数値根判定はこの既存構造の直接の帰結として証明し，
射影平面・二次曲面の格子分類は\cite[Theorem 5.12]{Kuz}に含まれる既知分類として扱う．

圏論的認識についてはMori--Ueyama\cite[Theorem 4.1]{MU}が，canonical bimoduleと
ample pairのfull geometric relative helixを用いてAS-regularなtailsを特徴付ける．
同論文の最終稿v4を用いる．v1のTheorem 3.10の証明が訂正され主定理が改訂されたことが
著者のarXiv記録に明示されており，旧版の強い結論を引用しない．
本研究で未解決なのは，full geometric helixをより弱い内在的条件から\emph{導く}部分である．

Hanihara\cite[Definition 2.5, Definitions 3.1--3.2, Theorem 3.3]{Hanihara}のroot pairは，
逆双対化dg双加群の根に生成・半直交条件を加え，巡回不変性の下でCalabi--Yau completionを与える．
数値的unipotenceはこの巡回不変性を意味せず，根だけからroot pairの生成条件も得られない．
従ってこの構成を使うためにfull性を仮定するなら，認識問題のfull性を解いたことにはならない．

Loginov--Orlov\cite[Definitions 2.1, 2.3, 2.6; Proposition 2.7]{LO}の非可換Fano多様体は，
smooth proper dg圏，例外対象の全非負反標準軌道上の高次消滅，指定された導来切断同値を要求する．
minifoldにはさらに長さ$n+1$のfull exceptional collectionとgeometric helix，反Serre作用の最大unipotenceを課す．
本稿の候補heartからこれらの消滅・同値・full性が自動的に従うとはしない．
逆向きに，その定義から本稿の構造対象のGabriel-rank条件が得られるとも主張しない．

従って本稿の価値の中心は，定義や既知格子分類の言い換えそれ自体ではなく，
自己同型の混入を正確に分離し，数値的整除条件を実際の消滅・境界・生成へ結び付けることにある．
新規性が未確認の初等的帰結はその旨を明記し，圏の実現問題を数値計算で代用しない．

\subsection{正規化の定義}

$\mathcal A$を$\C$線形Ext-finite abel圏とし，$D^b(\mathcal A)$のSerre関手を$\SSer$とする．
$\tau=\SSer^{-1}[n]$はheart上exactと仮定する．Euler形式の共通radicalで割って
\[
 N(\mathcal A)_{\Q}=\bigl(K_0(\mathcal A)/\operatorname{rad}\chi\bigr)\otimes\Q
\]
と書く．Serre双対性と自己同値による不変性から左右radicalは一致する．
数値群全体の有限次元性は，一般の候補圏に無断で仮定しない．

\begin{definition}[数値的に正規化された根]\label{def:normalized}
$\mathcal O\in\mathcal A$を固定し，$(\mathcal O,\tau)$をAZ ample pairとする．
heart上のexact自己同値$s$と自然同型$s^r\simeq\tau$を考える．
ある$M>0$に対して$(s_*-1)^M=0$が$\Nnum$全体で成り立つとき，この根を
\emph{数値的に正規化された根}と呼ぶ．
より弱い条件
\begin{equation}\label{eq:orbit-normalized}
 (s_*-1)^M[\mathcal O]=0\quad\text{in }\Nnum
\end{equation}
だけを課す場合は，\emph{構造対象の軌道上で正規化された根}と呼ぶ．
前者で可能な正整数$r$の集合を$\Runi(\mathcal A,\mathcal O)$とする．
これが非空有限の場合だけ，その最大値を$i_{\mathrm u}$，$n+1-i_{\mathrm u}$を正規化coindexと定める．
\end{definition}

$(\mathcal O,\tau)$がampleなら$(\mathcal O,s)$もampleである．
AZの生成条件には$s^{-rm}\mathcal O$だけを用いればよい．
全射に対するHomの漸近全射性は，指数を$rm+j$，$0\le j<r$に分け，
有限個の$s^j$で移した全射に$\tau$の条件を適用して従う．
この観察は数値正規化を必要としない．

\begin{remark}
正規化は構造対象の軌道だけでなく全数値群に課すのが本来の定義案である．
ただし以下の整除条件は\eqref{eq:orbit-normalized}だけで証明できる．
数値群全体のunipotence，軌道上のunipotence，Euler行列のrankを同一視しない．
また根の数値作用が存在しても，それをexact自己同値に持ち上げられるとは限らない．
\end{remark}

\section{古典的Fano指数との一致}
\begin{theorem}\label{thm:classical}
$X$を滑らかな連結射影複素Fano多様体とする．標準heartと$\mathcal O_X$について
\[
 \Runi(\coh X,\mathcal O_X)=\{r>0:r\mid i(X)\}.
\]
従って$i_{\mathrm u}(X)=i(X)$であり，正規化coindexは古典的coindexと一致する．
より詳しく，$s=L\otimes f^*(-)$と書いたとき，$s$が正規化された$r$乗根であることは
\[
 f^r=\id_X,\qquad L^{\otimes r}\simeq\omega_X^{-1},\qquad
 f^*|_{N(X)_{\Q}}=\id
\]
と同値である．ここで自己同値は同型まで考える．
\end{theorem}
\begin{proof}
Bondal--Orlov\cite[Theorem 3.1]{BO}により，導来自己同値は自己同型，線束によるtensor，shiftの合成である．
heartを保つのでshiftは零であり，$s=L\otimes f^*$と書ける．
$s^r=\tau$から，点対象への作用と$\mathcal O_X$への作用をそれぞれ見ると
\begin{equation}\label{eq:norm}
 f^r=\id_X,\qquad
 L\otimes f^*L\otimes\cdots\otimes(f^{r-1})^*L\simeq\omega_X^{-1}
\end{equation}
を得る．
数値$K$群のrank零部分から$c_1$によって$N^1(X)_{\Q}$への全射を得る．
これは数値的に定義される：点対象とのEuler pairingはrankを，曲線の構造層とのpairingはrank零の類の$c_1$の交点を検出する．
線束tensorはこの商で恒等であり，$s_*$の誘導作用は$f^*$となる．
従って$s_*$のunipotenceと$(f^*)^r=1$から，$f^*=1$ on $N^1(X)_{\Q}$である．
式\eqref{eq:norm}は$-K_X\equiv r c_1(L)$を与える．

ここでFano多様体では数値的自明な線束は自明であることを直接確認する．
$M\equiv0$なら$M\otimes\omega_X^{-1}$はampleなのでKodaira消滅より$H^{>0}(X,M)=0$．
数値Riemann--Rochから$\chi(M)=\chi(\mathcal O_X)=1$であり，$h^0(M)=1$となる．
その非零切断の有効divisorはample divisorとの交点が零なので零である．従って$M\simeq\mathcal O_X$．
よって$f^*L\simeq L$および$L^{\otimes r}\simeq\omega_X^{-1}$を得る．
Picard群は$H^2(X,\Z)$に同型で有限生成，かつ今の議論でtorsion-freeである．
従って$r\mid i(X)$である．

$f^*L=L$なので，数値作用において線束tensor $P$と$f^*$は可換する．
$P$はunipotentである．実際Riemann--Rochにより，全ての$F,G$について$\chi(F,L^{\otimes k}\otimes G)$は$k$の次数$\dim X$以下の多項式なので，
$(P-I)^{\dim X+1}=0$が数値群上成立する．
$f^*$は有限位数，従って標数零でsemisimpleである．
可換なunipotent作用とsemisimple作用の積$P f^*=s_*$がunipotentであるため，$f^*=1$が必要となる．
逆に表示された三条件から，norm条件とunipotenceが従う．
最後に$r\mid i(X)$なら$L^{\otimes r}=\omega_X^{-1}$となるampleな$L$を選び，
$s=L\otimes-$とすれば正規化された根を得る．
\end{proof}

\begin{proposition}[軌道上の正規化の正確な強さ]\label{prop:orbit-classical}
定理\ref{thm:classical}の$X$について，$s=L\otimes f^*$が$r$乗根であるとする．
次の三条件は同値である：
\[
 (s_*-I)^M[\mathcal O_X]=0\text{ for some }M,
 \qquad f^*L\simeq L,
 \qquad L^{\otimes r}\simeq\omega_X^{-1}.
\]
従って軌道上の正規化でも，許される次数全体は$i(X)$の約数に一致する．
ただし固定した根に対して，全数値群上の正規化はさらに$f^*|_{N(X)_{\Q}}=I$を要求する．
\end{proposition}
\begin{proof}
$T=s_*$，$e=[\mathcal O_X]$，$\ell=c_1(L)$とする．軌道上のunipotenceから$T^ke$は$k$のベクトル値多項式である．
rank零部分から$N^1(X)_{\Q}$への$c_1$を適用すれば，
\[
 c_1(s^k\mathcal O_X)=\sum_{j=0}^{k-1}(f^*)^j\ell
\]
も$k\ge0$で多項式である．その一次差分$(f^*)^k\ell$は$f^r=1$により周期的なので定数であり，$f^*\ell=\ell$となる．
数値的に自明な線束の自明性から$f^*L=L$．norm条件により$L^r=\omega_X^{-1}$を得る．
逆に$L^r=\omega_X^{-1}$なら$\operatorname{Pic}(X)$にtorsionがないので$f^*L=L$である．
このとき$s^k\mathcal O_X=L^k$となり，線束tensorの数値作用のunipotenceが軌道上の条件を与える．
\end{proof}

\begin{example}[二つの正規化が異なる根]\label{ex:weak-strong}
$X=\mathbb P^1\times\mathbb P^1$，$f$を因子交換とし，
$s=\mathcal O(1,1)\otimes f^*$とすると$s^2=\tau$である．
$s^k\mathcal O_X=\mathcal O(k,k)$なので軌道上では正規化されるが，
rulingの差には固有値$-1$を持つので全数値群では正規化されない．
さらに構造対象だけを$M=\mathcal O(1,0)$に変更すると
$c_1(s^kM)=(k,k)+c_1(f^{k*}M)$に周期成分が生じる．
従って弱い条件は構造対象に依存する．全体の条件を標識非依存の定義案として保持し，
弱い条件は個別の判定定理の必要仮定として用いる．
\end{example}

\begin{corollary}
前定理の設定で正規化された$(n+1)$乗根が存在すれば$X\simeq\mathbb P^n$である．
\end{corollary}
\begin{proof}
$-K_X\equiv(n+1)c_1(L)$かつ$L$はampleであるから，\cite[Corollary 2(1)]{KK}を適用する．
\end{proof}

\begin{example}[正規化前の根とnorm条件]\label{ex:non-normalized}
$X=\operatorname{Bl}_{p_1,p_2,p_3}\mathbb P^2$で，三つの座標点を巡回置換する$f$を取る．
$L=\mathcal O_X(H-E_1)$なら，$s=L\otimes f^*$は
$s^3=(-)\otimes\mathcal O_X(3H-E_1-E_2-E_3)=\tau$を満たす．
古典的指数は$1$だが三乗根が存在する．$f^*$は例外曲線の差の空間で非自明な三乗根の固有値を持つため正規化されない．
また$X=\mathbb P^1\times\mathbb P^1$で因子交換$f$と$L=\mathcal O(1,0)$を取ると
$s^2=(-)\otimes\mathcal O(1,1)$，$s^4=\tau$となる．
古典的指数は$2$であり，rulingの差に対する固有値$-1$が正規化を妨げる．
いずれも$s$はAZ ampleである．
\end{example}

\begin{remark}[残すべき自己同型]
正規化は$f=\id_X$を要求しない．例えば$\mathbb P^2$の位数3の自己同型$\alpha$は数値$K$群上恒等であり，
$\mathcal O(1)\otimes\alpha^*$は正規化された三乗根である．
その軌道環は非可換なZhang twistになり得る．この自然な非可換表示を正規化によって除外しないことが重要である．
\end{remark}

\begin{example}[naiveなcoindexのずれは任意に大きくなる]\label{ex:product}
$X=(\mathbb P^a)^m$で$m$因子を巡回置換する$f$と
$L=\mathcal O(1,0,\ldots,0)$を取る．$s=L\otimes f^*$は
\[
 s^m=(-)\otimes\mathcal O(1,\ldots,1),\qquad
 s^{m(a+1)}=\tau
\]
を満たす．$\dim X=ma$，通常の指数は$a+1$であり，根の次数だけから計算する値は
$\dim X+1-m(a+1)=1-m$となる．これは通常のcoindexではない．
自己同値のample性は$s^m$のample性から従う．
\end{example}

\begin{proposition}[標識変更と巡回群コホモロジー]\label{prop:root-cohomology}
滑らかな複素Fano多様体$X$，$f^r=\id_X$を固定し，$P=\operatorname{Pic}(X)$を加法記法で書く．
$F=f^*$，$N_r=I+F+\cdots+F^{r-1}$とする．
根$s=L\otimes f^*$の同型類は$N_r\ell=-K_X$の解に対応する．
解が一つ存在すれば，線束tensorによる共役を法とした解集合は
\[
 H^1(C_r,P)=\ker N_r/(F-I)P
\]
のtorsorである．ここでは同型$s^r\simeq\tau$の選択自体は標識しない．
$r\mid i(X)$なら唯一の$H\in P$で$rH=-K_X$となるものを基点にでき，
与えられた根が線束共役で軌道上正規化されるための必要十分条件は
\[
 [\ell-H]=0\quad\text{in }H^1(C_r,P)
\]
である．この群は有限で，$r$により消える．
\end{proposition}
\begin{proof}
式\eqref{eq:intro-norm}が第一の主張を与える．$M$によるtensor共役は
\[
 L\otimes f^*\longmapsto (L\otimes M\otimes f^*M^{-1})\otimes f^*
\]
なので$\ell\mapsto\ell+(I-F)m$である．解の差は$\ker N_r$に属する．
巡回群コホモロジーのこの表示は標準公式\cite[II, Example 1.20]{Milne}であるが，上の計算だけで同値関係も確認できる．
$rH=-K_X$なら$FH=H$であり，命題\ref{prop:orbit-classical}から軌道上正規化された線束は$H$だけである．
最後に$N_rx=0$なら
$(I-F)\sum_{j=0}^{r-1}jF^jx=-rx$だから商は$r$で消え，有限生成性から有限である．
\end{proof}

\begin{remark}
この命題は自己同型$f$を固定した線束共役の分類であり，$\operatorname{Aut}(X)$全体の共役分類ではない．
また，数値的に同じ根と互換的なdg双加群の根の構成を混同しない．
上の$H^1$は線束共役の障害であり，Haniharaの巡回不変性の障害と同じだとは主張しない．
\end{remark}

\section{正規化された根の数値作用とEuler多項式}
\begin{lemma}[一意な数値的根]\label{lem:unique-root}
$U$を標数零の有限次元ベクトル空間上のunipotent作用とする．
$T^r=U$を満たすunipotent作用は
\[
 T=\exp\bigl(r^{-1}\log U\bigr)
\]
だけである．特に$(U-I)^3=0$なら
\[
 T=I+\frac{U-I}{r}+\frac{1-r}{2r^2}(U-I)^2.
\]
この有理行列が指定された整数格子を保つことは，その格子上の正規化された根の存在の必要条件である．
\end{lemma}
\begin{proof}
unipotent作用のlogarithmとnilpotent作用のexponentialは有限和で互いに逆である．
$\log(T^r)=r\log T$より一意性と存在が従う．最後の式は二次までの二項展開である．
整数格子の保存は圏論的自己同値が存在するための必要条件にすぎず，十分条件とは主張しない．
\end{proof}

\begin{corollary}[有限位数部分の分離]\label{cor:finite-part}
標数零の有限次元数値空間上で$U$がunipotentとする．$T^r=U$を満たす全ての根は
\[
 T=U^{1/r}Q,\qquad Q^r=I,\qquad QU=UQ
\]
と一意に書ける．$U^{1/r}$は前補題のunipotent根であり，$Q$はsemisimpleである．
全体の正規化は$Q=I$，構造対象$e$の軌道上の正規化は$Qe=e$と同値である．
\end{corollary}
\begin{proof}
$T$は$T^r=U$と可換するので，$U$の多項式である$U^{1/r}$とも可換する．
$Q=U^{-1/r}T$と置けば$Q^r=I$であり，標数零なので$Q$はsemisimpleである．
$Q$の固有空間ごとに$T$は固有値をその固有値とするunipotent作用との積になる．
従って$T$がunipotentであるのは$Q=I$のときだけで，
$(T-I)^Me=0$であるのは$e$が$Q$の固有値1の空間に属するときだけである．
\end{proof}

\begin{proposition}[軌道の数値的生成性と正規化]\label{prop:orbit-spanning-normalization}
$N=K_0(\mathcal A)_{\Q}/\operatorname{rad}\chi$を有限次元とし，
$T=s_*$，$e=[\mathcal O]$とする．軌道の有理線形包
$V_e=\operatorname{span}_{\Q}\{T^ke\mid k\in\Z\}\subset N$について
\[
 (T-I)^Me=0\quad\Longleftrightarrow\quad
 (T|_{V_e}-I)^M=0
\]
が同じ$M>0$について成り立つ．特に$V_e=N$なら，
構造対象の軌道上の正規化は全数値群上の正規化と同値である．
\end{proposition}
\begin{proof}
$(T-I)^M$は$T$と可換するので，全ての整数$k$について
$(T-I)^MT^ke=T^k(T-I)^Me=0$である．従って$V_e$全体で零になる．
逆は$e\in V_e$から直ちに従う．
この結論はSerre関手やEuler形式のsurface-like性を仮定しない．
\end{proof}

\begin{remark}
例えば軌道対象の数値類が有限本で$N$を張ることが分かれば，
弱い正規化条件の標識依存性はこの範囲では消える．
ただし$V_e=N$を単なる反標準Hilbert関数，あるいはphantom残余を含むかもしれない
例外列の数値情報だけから仮定なしに導くことはできない．
\end{remark}


$T$が整数格子を保っても，$U^{1/r}$と$Q$が別々にその格子を保つとは限らない．
また，数値的有限位数部分$Q$を圏の自己同値として取り出せるとは限らない．
従ってこの分解だけで任意の圏論的根を正規化できるとは主張しない．

\begin{lemma}[軌道上の多項式性]\label{lem:polynomial}
$e=[\mathcal O]$とし，$T=s_*$が$(T-I)^Me=0$を満たすとする．
ある$F\in\Q[x]$について$\chi(\mathcal O,\tau^m\mathcal O)=F(m)$が$m\gg0$で成立すれば，
全整数$k$について
\[
 q(k):=\chi(\mathcal O,s^k\mathcal O)=F(k/r)
\]
である．
\end{lemma}
\begin{proof}
任意の正負の整数$k$について
\[
 q(k)=\sum_{j=0}^{M-1}\binom{k}{j}\chi(e,(T-I)^je)
\]
であり，これは$k$の有理多項式である．$q(rm)=F(m)$が無限個の$m$で成立するので多項式恒等式となる．
この証明に数値群全体の有限rankは不要である．
\end{proof}

\begin{theorem}[曲面指数の整除条件]\label{thm:divisibility}
$s^r=\tau=\SSer^{-1}[2]$を満たすexactな根が構造対象の軌道上で正規化され，
$d>0$について
\begin{equation}\label{eq:ac-euler}
 \chi(\mathcal O,\tau^m\mathcal O)=1+\frac d2m(m+1)\qquad(m\gg0)
\end{equation}
が成立するとする．このとき全整数$k$で
\begin{equation}\label{eq:surface-polynomial}
 q(k)=1+\frac{d}{2r^2}k(k+r).
\end{equation}
さらに
\begin{equation}\label{eq:integrality}
 a:=\frac d{r^2}\in\Z_{>0},\qquad a(r+1)\in2\Z
\end{equation}
である．すなわち$r^2\mid d$であり，$r$が偶数なら$2r^2\mid d$である．
\end{theorem}
\begin{proof}
前補題から\eqref{eq:surface-polynomial}を得る．整数値Euler列の二次差分は$a$なので$a\in\Z$であり，
$q(1)-q(0)=a(r+1)/2\in\Z$である．逆に\eqref{eq:integrality}は
\[
 q(k)=1+a\binom{k}{2}+\frac{a(r+1)}2k
\]
を全整数で整数値にする．この逆向きは多項式の整数値性についてだけで，圏や根の存在を与えない．
\end{proof}

\begin{corollary}[次数1から9の範囲]\label{cor:small-degree}
前定理で$1\le d\le9$なら，$r>1$の可能性は
\[
 (d,r)=(8,2),\quad(9,3)
\]
だけである．従って，この範囲では正規化された三乗根があれば$d=9$である．
\end{corollary}
\begin{proof}
$r^2\le d\le9$より$r=2,3$．$r=2$なら$8\mid d$，$r=3$なら$9\mid d$である．
\end{proof}

\begin{remark}[Hilbert関数とEuler多項式を区別する]\label{rem:vanishing}
反標準Hilbert関数$\dim\Hom(\mathcal O,\tau^m\mathcal O)=1+d m(m+1)/2$だけでは
\eqref{eq:ac-euler}は得られない．例えば
\[
 \Ext^i(\mathcal O,\tau^m\mathcal O)=0\quad(i>0,\ m\gg0)
\]
を別に確保すれば十分である．全対象の漸近Serre消滅はこれより強い十分条件である．
また$d\le9$は本稿では追加条件であり，一般の非可換Fano曲面について導いた上界ではない．
次数8での$r=2$は数値的必要条件であり，曲面や圏の認識を意味しない．
\end{remark}

\section{曲面型擬格子上の根と標準類}\label{sec:pseudolattice}
この節はKuznetsovの既存の数値的枠組み\cite[\S3--5]{Kuz}の中での，正規化根の判定である．
圏，heart，ample pair，Ext消滅を仮定した定理ではなく，それらを構成する結論も含まない．

\subsection{数値的線束tensorの明示式}
$G$を有限rankの自由整数群，$\chi$を非退化なunimodular整数双線形形式とする．
primitiveな$p\in G$が
\[
 \chi(p,p)=0,\quad\chi(p,v)=\chi(v,p),\quad
 \chi|_{p^\perp}\text{ は対称}
\]
を満たすとき，$(G,\chi,p)$はsurface-like pseudolatticeである\cite[Definition 3.1]{Kuz}．
$\rho(v)=\chi(p,v)$，$\operatorname{NS}(G)=p^\perp/\Z p$と置き，積$D\cdot E$を$-\chi$で定める．
標準類$K\in\operatorname{NS}(G)$は
\[
 \chi(v,w)-\chi(w,v)=-K\cdot\overline{\rho(v)w-\rho(w)v}
\]
によって定まる（バーは$\operatorname{NS}(G)=p^\perp/\Z p$への像を表す）．
存在と整性は\cite[Lemma 3.12]{Kuz}による．
以下，$K$がcharacteristic，すなわち$D^2\equiv K\cdot D\pmod2$であり，
$\rho(e)=1,\ \chi(e,e)=1$となる$e\in G$があるとする．

各$D\in\operatorname{NS}(G)$に対して，$\chi(e,x_D)=0$となる唯一のlift $x_D\in p^\perp$を選ぶ．
$\chi(e,p)=1$なのでこれは加法的に選べ，$G=\Z e\oplus x_{\operatorname{NS}(G)}\oplus\Z p$となる．
座標$v=(a,D,c)=ae+x_D+cp$でEuler形式は
\begin{equation}\label{eq:lattice-euler}
 \chi((a,D,c),(b,E,d))=ab+ad+bc+bK\cdot D-D\cdot E.
\end{equation}

\begin{lemma}[数値的tensor作用]\label{lem:translation}
$H\in\operatorname{NS}(G)$について
\begin{equation}\label{eq:numerical-tensor}
 t_H(a,D,c)=\left(a,D+aH,c+D\cdot H+\frac a2(H^2-K\cdot H)\right)
\end{equation}
は整数格子$G$のunipotent isometryである．さらに
$t_Ht_J=t_{H+J}$，$t_H^{-1}=t_{-H}$，Serre作用は$S=t_K$である．
\end{lemma}
\begin{proof}
characteristic条件により最後の成分は整数である．
式\eqref{eq:lattice-euler}へ代入すればisometryが従い，二次式の交差項$H\cdot J$から合成公式を得る．
$t_H-I$はfiltration $0\subset\Z p\subset p^\perp\subset G$を一段下げるので三乗は零である．
また式\eqref{eq:lattice-euler}から$\chi(v,w)=\chi(w,t_Kv)$であり，非退化性から$S=t_K$となる．
\end{proof}

\begin{theorem}[unipotentな数値根の完全判定]\label{thm:lattice-root}
上の仮定の下で，unipotentな整数格子自己同型$T$で$T^r=S^{-1}$を満たすものが存在するための必要十分条件は
\[
 \boxed{-K\in r\operatorname{NS}(G).}
\]
存在するなら$H=-K/r$として$T=t_H$だけであり，自動的に$\chi$を保つ．
従って$K\ne0$なら，可能な次数は$\operatorname{NS}(G)$における標準類のdivisibilityの約数全体である．
\end{theorem}
\begin{proof}
有理係数では補題\ref{lem:unique-root}からunipotent根は一意である．
式\eqref{eq:numerical-tensor}を有理$H$まで延長すれば$t_{-K/r}^r=t_{-K}=S^{-1}$なので，これが唯一の根である．
この作用が$G$を保つなら$T(e)$の$\operatorname{NS}$成分$-K/r$は整数である．
逆に$H=-K/r$が整数なら補題\ref{lem:translation}が必要な整数isometryを与える．
これは\emph{数値的自己同型}の存在の十分条件であり，圏の自己同値の存在の十分条件ではない．
\end{proof}

\begin{corollary}\label{cor:lattice-euler}
$d=K^2>0$，$T^r=S^{-1}$を満たすunipotent整数根$T$について
\[
 \chi(e,T^ke)=1+\frac{d}{2r^2}k(k+r).
\]
従って$r^2\mid d$，偶数$r$なら$2r^2\mid d$である．
この節の仮定では，この結論に漸近Ext消滅も反標準Hilbert関数も必要ない．
代わりに，point-like類，rank-one類，characteristic標準類を持つ指定した整数擬格子を仮定している．
\end{corollary}
\begin{proof}
$T^k=t_{kH}$を式\eqref{eq:numerical-tensor}へ代入する．$d=r^2H^2$と
$H^2\equiv K\cdot H=-rH^2\pmod2$が整除条件を与える．
\end{proof}

\subsection{正規化しない根の半直積表示と周期的Euler欠損}
前小節はunipotentな根を扱った．ここでは数値的unipotenceを先に課さず，
有限位数の作用が何を変えるかを明示する．次の表示は曲面型擬格子上の
数値等長変換の初等的な計算であり，一般のabel圏の自己同値群を分類するものではない．

$\operatorname{O}(\operatorname{NS}(G))_K$を，交点形式を保ち$K$を固定する
整数等長変換の群とする．$F$がこの群に属するとき，
$\widetilde F(a,D,c)=(a,FD,c)$と定める．

\begin{theorem}[正規化しない数値根の分類]\label{thm:twisted-lattice-root}
前小節のunimodular surface-like擬格子，primitiveな$p$，rank-one exceptional類$e$，
characteristic標準類$K$の仮定を置く．$S=t_K$をSerre作用とする．
\begin{enumerate}[label=(\roman*)]
\item $p$を固定するEuler等長変換は，一意的に
\[
 T=t_H\widetilde F,\qquad H\in\operatorname{NS}(G),\quad
 F\in\operatorname{O}(\operatorname{NS}(G))_K
\]
と書ける．合成法則は
$t_H\widetilde F\,t_J\widetilde G=t_{H+FJ}\widetilde{FG}$である．
\item この形の$T$が$T^r=S^{-1}$を満たす必要十分条件は
\[
 F^r=I,\qquad N_r(F)H=-K,\quad
 N_r(F)=I+F+\cdots+F^{r-1}
\]
である．特に$K^2=d>0$なら$r\mid d$である．
全数値群上のunipotenceは$F=I$と同値で，軌道上の条件
$(T-I)^Me=0$（ある$M>0$）は$FH=H$，すなわち$H=-K/r$と同値である．
\item $K^2\ne0$とする．$T^r=S^{-1}$を満たす整数Euler等長変換で
$\rho(Te)>0$となるものは，自動的に$p$を固定するので(i)--(ii)が適用できる．
特に，数値rankがheartの対象に非負で，$e$がrank oneの構造対象ならば，
heart-exactな根はこの場合に属する．
\end{enumerate}
\end{theorem}
\begin{proof}
$p$を固定する$T$は$\rho(v)=\chi(p,v)$を保存する．
前小節の座標で$Te=(1,H,c)$，$T(0,D,0)=(0,FD,\lambda_D)$と書ける．
Euler形式の保存から$F$はNSの等長変換であり，交代化したEuler形式から$FK=K$が従う．
$\chi(Te,Te)=1$，$\chi(Te,T(0,D,0))=\chi(e,(0,D,0))=0$を順に用いると
\[
 c=\frac{H^2-K\cdot H}{2},\qquad
 \lambda_D=H\cdot FD
\]
となる．$Tp=p$と合わせて，これは正確に$T=t_H\widetilde F$である．
逆に補題\ref{lem:translation}と$FK=K$から$t_H\widetilde F$はEuler等長変換である．
共役式$\widetilde F t_H\widetilde F^{-1}=t_{FH}$が合成法則を与える．

従って$T^r=t_{N_r(F)H}\widetilde F^r$となり，(ii)の根の条件が従う．
$F^r=I$かつ$FK=K$なので，$K$との交点を取れば
$r(K\cdot H)=-K^2=-d$である．これより$r\mid d$．
$F$は有限位数だから，そのNS上の作用がunipotentであるのは$F=I$のときに限る．
$F=I$なら$T=t_H$はunipotentである．

$M_k=\sum_{j=0}^{k-1}F^jH$とすると，$k\ge0$について
$T^ke=t_{M_k}e$である．$(T-I)^Me=0$ならこの軌道は$k$のベクトル値多項式になり，
NS成分$M_k$も多項式になる．しかし差分$M_{k+1}-M_k=F^kH$は周期的なので一定であり，
$FH=H$である．逆に$FH=H$ならnorm条件から$rH=-K$で，
$T^ke=t_{kH}e$は$k$の二次多項式となる．

最後に$U=S^{-1}=t_{-K}$と置く．$U-I$は三段のrank filtrationを一段ずつ下げ，
$(U-I)^2e=K^2p\ne0$だから$\operatorname{im}(U-I)^2=\Q p$である．
$T^r=U$なら$TU=UT$なので$T$はこの一次元空間を保つ．
$p$はprimitiveな整数類であるから$Tp=\pm p$．
$Tp=-p$なら$\rho(Te)=\chi(p,Te)=\chi(T^{-1}p,e)=-1$となり
仮定に反する．よって$Tp=p$である．
\end{proof}

\begin{corollary}[周期的なEuler欠損と正規化の検出]\label{cor:periodic-deficit}
定理\ref{thm:twisted-lattice-root}に加え，
$\operatorname{NS}(G)$のsignatureを$(1,\operatorname{rk}G-3)$，$d=K^2>0$とする．
根$T=t_H\widetilde F$に対して
\[
 v=H+K/r\in K^\perp_{\Q},\qquad
 v_k=\sum_{j=0}^{k-1}F^jv\quad(k\ge0)
\]
とおく（負の$k$へは和の差分関係で延長する）．このとき$N_r(F)v=0$であり，
$v_{k+r}=v_k$である．全整数$k$について
\begin{equation}\label{eq:periodic-deficit}
 \boxed{\quad\chi(e,T^ke)=1+\frac{d}{2r^2}k(k+r)+\frac12v_k^2.\quad}
\end{equation}
$K^\perp_{\Q}$は負定値なので，最後の項は非正で，$r$を周期とし，$k\equiv0\pmod r$で零になる．
特に次は同値である：
\[
 (T-I)^Me=0\text{ for some }M,\quad
 v=0,\quad
 \chi(e,Te)=1+\frac{d(r+1)}{2r^2}.
\]
従って，正規化しない根のEuler軌道は，正規化した根の仮想二次多項式から
非正の周期的誤差だけずれる．
\end{corollary}
\begin{proof}
$N_r(F)H=-K$と$FK=K$の交点から$K\cdot H=-d/r$となり，$K\cdot v=0$．
また$N_r(F)v=N_r(F)H+K=0$である．
$M_k=\sum_{j=0}^{k-1}F^jH=-kK/r+v_k$，$M_{k+r}=M_k-K$．
式\eqref{eq:numerical-tensor}から
$\chi(e,T^ke)=1+\bigl(M_k^2-K\cdot M_k\bigr)/2$なので，
$K\cdot v_k=0$を使えば表示式が従う．$K^2>0$とsignatureの仮定から
$K^\perp$は負定値であり，$v_k^2=0$なら$v_k=0$．
最後の同値は$k=1$で$v_1=v$，および前定理(ii)による．
\end{proof}

\begin{corollary}[数値rank三では正規化が自動的]\label{cor:rank-three-normalized}
定理\ref{thm:twisted-lattice-root}の仮定で$K^2>0$かつ
$\operatorname{rk}G=3$とする．$\rho(Te)>0$となる根$T$はすべてunipotentである．
従ってその次数は$-K\in r\operatorname{NS}(G)$を満たすものに限られる．
これは数値rank三と曲面型格子の仮定の下での結論であり，
元の六公理からこれらの格子仮定を導いたという主張ではない．
\end{corollary}
\begin{proof}
$\operatorname{NS}(G)$はrank oneで，その非零元$K$を固定する等長変換は恒等変換だけである．
ゆえに前定理の$F=I$となる．
\end{proof}

\begin{proposition}[固定した数値自己同型と巡回障害]\label{prop:numerical-root-cohomology}
$F\in\operatorname{O}(\operatorname{NS}(G))_K$，$F^r=I$を固定する．
$N_r(F)H=-K$に整数解があれば，解を数値的線束移動$t_A$による共役で割った集合は
$H^1(C_r,\operatorname{NS}(G))=\ker N_r(F)/(I-F)\operatorname{NS}(G)$のtorsorである．
この群は有限で$r$により消える．有理解としては常に$-K/r$があり，
任意の整数解は$\operatorname{NS}(G)_{\Q}$上の数値移動によりそれと共役である．
ただし有理移動は通常，整数格子の自己同型でも圏の自己同値でもない．
\end{proposition}
\begin{proof}
共役で$H$は$H+(I-F)A$に変わる．二解の差は$\ker N_r(F)$であり，
標準的な巡回群コホモロジーの表示が従う．
$(I-F)\sum_{j=0}^{r-1}jF^j=N_r(F)-rI$なので群は$r$-torsionで，有限生成性から有限である．
有理係数では有限群の一次コホモロジーは零であるから，
$H+K/r\in\ker N_r(F)$は$(I-F)\operatorname{NS}(G)_{\Q}$に属する．
\end{proof}

\begin{example}[次数6のdel Pezzo曲面にある六乗根]\label{ex:hexagon-sixth-root}
$X=\operatorname{Bl}_{p_1,p_2,p_3}\mathbb P^2$を三つの座標点のblowupとする．
$\operatorname{Pic}(X)=\Z H\oplus\Z E_1\oplus\Z E_2\oplus\Z E_3$，
$K=-3H+E_1+E_2+E_3$で$K^2=6$である．
六本のtoric境界曲線を巡回順に
\[
 D_0=E_1,\ D_1=H-E_1-E_2,\ D_2=E_2,\
 D_3=H-E_2-E_3,\ D_4=E_3,\ D_5=H-E_3-E_1
\]
と書くと$D_0+\cdots+D_5=-K$である．
この曲面の六角形fanは$\begin{psmallmatrix}0&-1\\1&1\end{psmallmatrix}\in
\operatorname{GL}_2(\Z)$で一つずつ回転されるので，位数6のtoric自己同型$f$があり，
向きを選べば$f^*D_j=D_{j+1}$（添字は$6$を法とする）となる．
$s=\mathcal O_X(D_0)\otimes f^*(-)$とすると
\[
 s^6=(-)\otimes\mathcal O_X\Bigl(\sum_{j=0}^{5}D_j\Bigr)
 =(-)\otimes\omega_X^{-1}=\tau.
\]
従って$s$はheart-exactで，$s^6$が反標準線束のtensorなのでAZ ampleである．
古典的Fano指数は$1$であるにもかかわらず，正規化しない圏論的根には次数$6$がある．
一方，すべてのheart-exact根の次数は定理\ref{thm:twisted-lattice-root}により$6$を割るので，
この曲面の正規化しないexact指数は正確に$6$である．
しかも定理\ref{thm:twisted-lattice-root}の$r\mid K^2$はこの例で等号となる．
$r=6$，$e=[\mathcal O_X]$について$\chi(e,se)=1$だが，仮想的な正規化多項式の値は$19/12$であり，
式\eqref{eq:periodic-deficit}の欠損は$-7/12$である．
これは整数格子上の数値模型にとどまらず，実在する滑らかなFano曲面上のexact根である．
\end{example}

\subsection{defectによる次数上界と既知の格子分類}
$\operatorname{NS}(G)$のsignatureを$(1,\operatorname{rk}G-3)$とする．
この条件とcharacteristic性がKuznetsovのgeometricityである．defectは既知の定義
\begin{equation}\label{eq:defect}
 \delta(G)=K^2+\operatorname{rk}G-12
\end{equation}
を用いる\cite[Definition 5.3]{Kuz}．通常の$q=p_g=0$の曲面ではNoether公式から零になるが，
一般の非可換候補圏で零あるいは非正になることはまだ仮定である．

\begin{corollary}[非正defectでの数値的指数上界]\label{cor:defect-index}
この節の仮定に加えて$K^2>0$，$\delta(G)\le0$とする．
$T^r=S^{-1}$となるunipotent整数根について$r\le3$であり，
\begin{itemize}
\item $r=3$なら$K^2=9$，$\operatorname{rk}G=3$，$\delta(G)=0$で，
$G$は$K_{\mathrm{num}}(\mathbb P^2)$の標準数値型である．
\item $r=2$なら$K^2=8$，$\operatorname{rk}G=4$，$\delta(G)=0$で，
$G$は$K_{\mathrm{num}}(\mathbb P^1\times\mathbb P^1)$の標準数値型である．
\end{itemize}
後半の二つの格子分類は\cite[Theorem 5.12]{Kuz}に含まれる．ここでは根の条件からその仮定へ至る過程を直接記す．
\end{corollary}
\begin{proof}
$K^2>0$より$\operatorname{rk}G\ge3$．式\eqref{eq:defect}から$0<K^2\le9$であり，前系の整除条件により非自明な場合は$(d,r)=(9,3),(8,2)$だけである．
$r=3$では$\operatorname{rk}G\le3$なので等号で，$H=-K/3$の平方は1となる．
$e,Te,T^2e$のEuler行列は
\[
 \begin{pmatrix}1&3&6\\0&1&3\\0&0&1\end{pmatrix}
\]
で行列式1である．全格子のrankは3なので，この三類は整数基底をなし射影平面のEuler格子を与える．

$r=2$では$d=8$，$\operatorname{rk}G\le4$．もしrankが3ならNSは正定値rank 1のunimodular格子であり，
$H=-K/2$の平方2は実現できない．従ってrankは4，defectは零である．
$K$は2で割れるcharacteristic類なのでNSはevenで，signature $(1,1)$のunimodular格子，すなわち双曲平面である．
双曲基底$H_1,H_2$では符号を調整して$H=H_1+H_2$となる．
$e,t_{H_1}e,t_{H_2}e,t_{H_1+H_2}e$のEuler行列は
\[
 \begin{pmatrix}1&2&2&4\\0&1&0&2\\0&0&1&2\\0&0&0&1\end{pmatrix}
\]
となり，二次曲面の格子を与える．
\end{proof}

\begin{remark}[例外基底がある場合と圏論的限界]\label{rem:defect-exceptional}
幾何的surface-like擬格子がexceptional basisを持てば，
Kuznetsovの数値的縮約とdefectの単調性，最小型のdefect零から$\delta(G)\le0$が従う
\cite[Corollary 5.7, Lemma 5.8, Theorem 5.9]{Kuz}．
従ってこの範囲ではdefectの非正性を別に仮定しなくてよい．
ただし本稿の一般候補圏の数値群がgeometric surface-likeであること，指定したpoint-like類とgeneric rankが整合すること，
またexceptional basisやfull exceptional collectionの存在は証明していない．
数値的な射影平面型からphantom型の残余圏の零性を結論することもできない．
\end{remark}

\section{滑らかなdel Pezzo曲面の非正規化exact指数}\label{sec:delpezzo-exact}

以下は\emph{数値的な可能性}だけでなく，実際の滑らかな複素Fano曲面上にある
heart-exact自己同値の根を扱う．根$s$自身に線束のample性を課す必要はない：
$s^r=\tau$で$\tau$が反標準ampleなら$(\OO_X,s)$はAZ ampleである．
以下の幾何学的構成に使う自己同型，del Pezzo分類，例外曲線の記述は
\cite{DolgachevMartin,BauerCatanese,Rieken}と比較する．
これらの古典的事実自体を新しい成果と主張しない．

\begin{definition}[非正規化exact指数]\label{def:exact-index-surface}
滑らかな複素Fano曲面$X$に対し
\[
 i_{\mathrm{ex}}(X)=\max\bigl\{r\ge1\mid
 \exists\,s\in\operatorname{Aut}_{\mathrm{exact}}(\coh X),\quad
 s^r\simeq(-)\otimes\omega_X^{-1}\bigr\}
\]
と定める．最大値の存在は次の命題から従う．自己同値は自然同型まで考える．
正規化指数$i_{\mathrm u}(X)$とは区別する．
\end{definition}

\begin{proposition}[反標準次数による普遍的上界]\label{prop:geometric-exact-degree-bound}
滑らかな複素Fano曲面$X$の反標準次数を$d=K_X^2$とする．
$s^r\simeq\omega_X^{-1}\otimes(-)$を満たすheart-exact自己同値が存在すれば
\[
 \boxed{\quad r\mid d.\quad}
\]
特に$i_{\mathrm{ex}}(X)\mid d$である．
\end{proposition}
\begin{proof}
Bondal--Orlov\cite[Theorem 3.1]{BO}によって$s=L\otimes f^*$と書ける．
命題\ref{thm:classical}の証明にあるnorm条件は
$f^r=1$と$\sum_{j=0}^{r-1}(f^j)^*c_1(L)=-K_X$である．
$f^*K_X=K_X$なので，$-K_X$との交点を取ると
\[
 d=(-K_X)^2
 =\sum_{j=0}^{r-1}(-K_X)\cdot (f^j)^*c_1(L)
 =r\,(-K_X)\cdot c_1(L)\in r\Z.
\]
数値Grothendieck群の全体または軌道上でのunipotenceは用いていない．
\end{proof}

\begin{lemma}[巡回する反標準因子]\label{lem:cyclic-anticanonical-boundary}
$X$を滑らかな複素Fano多様体，$f\in\operatorname{Aut}(X)$を位数$r$の自己同型とし，
Cartier因子$D$が
\[
 \sum_{j=0}^{r-1}(f^j)^*D\sim-K_X
\]
を満たすとする．このとき$s=\mathcal O_X(D)\otimes f^*$はheart-exactで
$s^r\simeq\omega_X^{-1}\otimes(-)$を満たし，$(\mathcal O_X,s)$はAZ ampleである．
特に，反標準因子を構成する$r$個の因子を$f$が巡回する場合に適用できる．
\end{lemma}
\begin{proof}
線束とpullbackの合成を$r$回繰り返すと
$s^r\simeq\mathcal O_X(\sum_j(f^j)^*D)\otimes(f^r)^*$である．
従って仮定から反標準tensorと同型になる．反標準線束はampleなので，
$s^r$に対するAZの生成条件と漸近全射条件を指数を$r$で割った剰余類に分ければ
$(\mathcal O_X,s)$のample性が従う（定義\ref{def:normalized}の直後の議論）．
\end{proof}

\begin{theorem}[del Pezzo次数ごとの実現可能な最大exact指数]
\label{thm:delpezzo-exact-index-table}
$d\in\{1,\ldots,9\}$について，$K_X^2=d$となる滑らかな複素del Pezzo曲面全体を
動かしたとき，$i_{\mathrm{ex}}(X)$の最大値は
\[
\begin{array}{c|ccccccccc}
 d&1&2&3&4&5&6&7&8&9\\\hline
 \max_{K_X^2=d} i_{\mathrm{ex}}(X)&1&2&3&4&5&6&1&4&3
\end{array}
\]
である．次数$2$では\emph{すべての}滑らかなdel Pezzo曲面で指数$2$，
次数$5,6$でも各次数の曲面は同型を除いて一つなので指数はそれぞれ$5,6$である．
次数$8$では$\mathbb F_1$の指数は$1$，$\mathbb P^1\times\mathbb P^1$の指数は$4$である．
次数$3,4$については最大値を達成する滑らかな曲面を具体的に構成するが，
各次数の\emph{任意の}曲面がその最大値を達成すると主張しない．
\end{theorem}
\begin{proof}
上界は命題\ref{prop:geometric-exact-degree-bound}である．次数$1$は直ちに指数$1$．
残る次数について構成と上界の鋭さを順に説明する．

\emph{次数$2$．}
反標準写像$\pi:X\to\mathbb P^2$は滑らかな四次曲線を分岐とする二重被覆であり，
そのdeck involutionを$\gamma$（Geiser involution）とする．
$X$の任意の$(-1)$-曲線$E$は$(-K_X)\cdot E=1$なので，
$\pi(E)$は直線$\ell$で，$E$は$\pi^{-1}(\ell)$の一成分である．
二重被覆のdeck交換により
$\pi^*\ell=E+\gamma(E)\sim-K_X$となる．
ゆえに補題\ref{lem:cyclic-anticanonical-boundary}で$r=2$を得る．
この構成は任意の滑らかな次数2曲面で成立する．

\emph{次数$3$．}
Fermat三次曲面
\[
 X=\{x^3+y^3+z^3+w^3=0\}\subset\mathbb P^3
\]
と原始三乗根$\zeta$を取る．$f(w)=\zeta w$という射影線形自己同型は位数$3$である．
$X$の超平面切断$\{x+y=0\}$は
\[
 z^3+w^3=\prod_{j=0}^2(z+\zeta^jw)=0
\]
なので，三直線$D_j=\{x+y=z+\zeta^jw=0\}$の和である．
$f$がこれらを巡回し，$D_0+D_1+D_2\sim-K_X$であるから三乗根を得る．

\emph{次数$4$．}
$Y=\mathbb P^1\times\mathbb P^1$で
$f(u,v)=(v,u^{-1})$とする．$f$は位数$4$であり，
$t\in\C^*$，$t^4\ne1$に対して
\[
 p_0=(0,t),\qquad p_1=(t,\infty),\qquad
 p_2=(\infty,t^{-1}),\qquad p_3=(t^{-1},0)
\]
を巡回する．四点の第一・第二座標はそれぞれ相異なる．
$(1,1)$形式$a uv+bu+cv+d$の四点での値を並べる行列の行列式は，
非零の座標倍を除けば$(t^4-1)/t^2$であり，ゼロでない．
従って同じruling fiber上の二点はなく，四点を通る$(1,1)$曲線もない．
四点は次数$4$ del Pezzo曲面のblowupに適した一般位置にあり，
$X=\operatorname{Bl}_{\{p_0,p_1,p_2,p_3\}}Y$は滑らかなdel Pezzo曲面である．
$f$は$X$へ持ち上がる．$Y$の四本のtoric境界成分のstrict transformを
$D_0,\ldots,D_3$とすると，$f$がこれらを巡回し，
各境界成分は一点でblowupされるので
\[
 \sum_{j=0}^3 D_j\sim \pi^*(-K_Y)-\sum_{j=0}^3E_j=-K_X.
\]
従って四乗根が存在する．例えば$t=2$を選べる．

\emph{次数$5$．}
この次数の滑らかな曲面$X$は同型を除いて一意であり，
$\operatorname{Aut}(X)\simeq\mathfrak S_5$である\cite{BauerCatanese}．
$\operatorname{Pic}(X)$の標準類の直交補は$A_4$の根格子で，
$\mathfrak S_5$は標準のWeyl作用を持つ．
位数$5$のcycle $f$は$A_4\otimes\Q$に不動ベクトルを持たないため，
$\operatorname{Pic}(X)_{\Q}^{\langle f\rangle}=\Q K_X$である．
任意の$(-1)$-曲線$E$のnorm
$N_5(E)=\sum_{j=0}^4(f^j)^*E$は$K_X$の整数倍であり，
$K_X\cdot N_5(E)=5K_X\cdot E=-5$，$K_X^2=5$だから
$N_5(E)=-K_X$である．補題により五乗根が存在する．

\emph{次数$6$．}
例\ref{ex:hexagon-sixth-root}の六角形fanの回転と反標準toric境界から
六乗根を得る．この次数の滑らかなdel Pezzo曲面は同型を除いて一つである．

\emph{次数$7$．}
これは$\mathbb P^2$の二点blowupであり，$(-1)$-曲線は
$E_1,E_2,H-E_1-E_2$の三本だけである．
最後の曲線は他の二本と交わり，残る二本は互いに交わらないので，
任意の自己同型のPicard群上の作用は恒等または$E_1,E_2$の交換である．
根の次数$r$は$7$を割るから$r>1$なら$r=7$である．
この場合$f^7=1$なのでPicard作用は恒等となり，norm条件は$7L=-K_X$となる．
しかし$-K_X$はprimitiveであり不可能である．従って指数は$1$．

\emph{次数$8$．}
滑らかな曲面は$\mathbb F_1$か$\mathbb P^1\times\mathbb P^1$である．
$\mathbb F_1$では唯一の$(-1)$-曲線と$K_X$が$\operatorname{Pic}(X)_{\Q}$を生成するので
全自己同型がPicard群上恒等であり，$-K_X$のprimitive性から指数は$1$である．
$\mathbb P^1\times\mathbb P^1$ではPicard群上の自己同型は恒等か二つのruling交換である．
$-K_X=(2,2)$なので$r=8$の根があると，normは恒等作用の場合$8(a,b)$，
交換の場合$4(a+b,a+b)$となり，いずれも$(2,2)$にならない．
他方，ruling交換$f$と$L=\mathcal O(1,0)$により
$(L\otimes f^*)^4=\tau$が成立するため指数は$4$である．

\emph{次数$9$．}
$X=\mathbb P^2$でPicard群上の自己同型作用は常に恒等である．
$-K_X=3H$なので根の次数は$3$を割り，$\mathcal O(1)\otimes-$が三乗根となる．
以上で表の全てが従う．
\end{proof}

\begin{proposition}[同じ数値格子でもexactな四乗根は変わる]
\label{prop:degree-four-lifting-gap}
滑らかな複素del Pezzo曲面$X$で$K_X^2=4$とする．
\begin{enumerate}[label=(\roman*)]
\item 全ての$X$で$i_{\mathrm{ex}}(X)\ge2$である．
\item 一般の$X$では$i_{\mathrm{ex}}(X)=2$である．
\item 一方，全ての$X$について，数値Euler格子$K_{\mathrm{num}}(X)$上には
反Serre作用$\tau_*$の\emph{整数格子を保つ四乗根}が存在する．
従って一般の$X$では，この数値根はheart-exact自己同値へ持ち上がらない．
\end{enumerate}
特に，数値Grothendieck格子の同型類から非正規化exact指数は決まらない．
\end{proposition}
\begin{proof}
次数$4$の曲面は二つの二次式の滑らかな完全交叉で，複素数体上では
\[
 X=\Bigl\{\sum_{i=0}^4x_i^2=\sum_{i=0}^4\lambda_i x_i^2=0\Bigr\}
 \subset\mathbb P^4,\qquad\lambda_i\ne\lambda_j\ (i\ne j)
\]
と同時対角化できる\cite{DolgachevMartin,HassettKreschTschinkel}．
$j$番目の特異二次曲面$Q_j=\{\sum_{i\ne j}(\lambda_i-\lambda_j)x_i^2=0\}$は
頂点を持つrank $4$のquadric coneであり，二組のplane rulingがある．
$X$に切って得る二組のconic pencilのfiber類を$A_j,B_j$とすると
$A_j+B_j=-K_X$である（rulingを跨ぐ二平面の和は超平面切断）．
$j\ne i$を選び，座標符号変更$f_i:x_i\mapsto-x_i$を取る．
これは$X$の位数2の自己同型で，$Q_j$の非特異quadric底面では行列式$-1$の直交反射なので
二組のrulingを交換する．従って$f_i^*A_j=B_j$となる．
任意のfiber $C$の因子類$[C]=A_j$を取れば
$C+f_i^*C\sim A_j+B_j=-K_X$なので補題\ref{lem:cyclic-anticanonical-boundary}から
exactな二乗根を得る．これで(i)を示した．

一般の次数$4$曲面の自己同型群が$(\Z/2)^4$であることは
\cite[Introduction]{HassettKreschTschinkel}の既知結果である．
従って一般の$X$で$f\in\operatorname{Aut}(X)$なら$f^2=1$であり，
$L\otimes f^*$の四乗の線束因子は
$2(L+f^*L)$（Picard群の加法記法）となる．
これは$2\operatorname{Pic}(X)$に属する一方，$-K_X$はprimitiveなので四乗根にならない．
命題\ref{prop:geometric-exact-degree-bound}から可能な指数は$1,2,4$に限られ，
(i)より一般の$X$では指数$2$である．

最後に，次数$4$の滑らかなdel Pezzo曲面はすべて$\mathbb P^2$の一般位置五点blowupである．
従って各$X$の$K_{\mathrm{num}}(X)$は，指定したrank-one類・point類・標準類を保ち，
同じ積形式とEuler形式を持つunimodular surface-like格子として同型である
（式\eqref{eq:lattice-euler}，またはRiemann--Roch）．
定理\ref{thm:delpezzo-exact-index-table}の特殊な四乗根が誘導する整数Euler等長変換を，
この同型で任意の$X$の数値格子へ移せば，$\tau_*$の四乗根が得られる．
一般の$X$では(ii)によりexactな四乗根が存在しないので，数値根の持上げは不可能である．
\end{proof}


\begin{remark}[この表の意味と新規性の範囲]\label{rem:delpezzo-exact-scope}
これは滑らかなdel Pezzo曲面の\emph{すべての自己同型を個別に分類した表}ではなく，
各次数における最大exact指数の実現可能値を決定した表である．
次数$3,4$の特別な実現を一般の同次数曲面に拡張していない．
正規化した指数は定理\ref{thm:classical}により古典的指数に戻り，
例えば次数$2,3,4,5,6$の例では正規化の有無で最大値が異なる．
構成に使ったFano曲面と自己同型の幾何学は既知であり，
これを圏論的exact指数として統一的に並べた結果の独立した新規性は未確認である．
非可換なAS-regular圏への根の持上げ，あるいは新しい非可換圏の認識はこの表からは従わない．
\end{remark}


\section{数値的Kobayashi--Ochiai型上界とその限界}
\begin{proposition}[負の窓を仮定した指数上界]\label{prop:window}
補題\ref{lem:polynomial}の設定で$F$の次数を$n$，$q(0)=1$とする．さらに
\[
 q(-j)=0\qquad(1\le j\le r-1)
\]
を仮定する．このとき$r\le n+1$であり，等号なら
\[
 q(k)=\binom{k+n}{n},\qquad F(m)=\binom{(n+1)m+n}{n}.
\]
\end{proposition}
\begin{proof}
$q$は零でない次数$n$の多項式で，相異なる$r-1$個の根を持つ．
等号の場合は根$-1,\ldots,-n$と$q(0)=1$が$q$を一意に決める．
\end{proof}

これは数値的な指数上界であり，空間や圏の認識ではない．
古典的Fano多様体では線束の負の窓の消滅にKodaira消滅を使えるが，
非可換の場合にはこの窓の消滅が独立した研究課題となる．
正規化だけでこの仮定を消すことはできないことを，次の整数格子模型が示す．

\begin{proposition}[任意に大きい奇数指数の数値模型]\label{prop:lattice-model}
任意の正の奇数$r$について，$N=\Z^3$，
\[
 T=\begin{pmatrix}0&0&1\\1&0&-3\\0&1&3\end{pmatrix},\qquad
 G_r=\begin{pmatrix}
 1&(r+3)/2&r+3\\
 (3-r)/2&1&(r+3)/2\\
 3-r&(3-r)/2&1
 \end{pmatrix}
\]
と置く．$\chi(x,y)=x^{\mathsf t}G_ry$，$e=(1,0,0)^{\mathsf t}$とすると
\[
 (T-I)^3=0,\quad\det G_r=1,\quad T^{\mathsf t}G_rT=G_r,
 \quad G_rT^{-r}=G_r^{\mathsf t},
\]
\[
 \chi(e,T^ke)=1+\tfrac12k(k+r),\qquad
 \chi(e,T^{rm}e)=1+\tfrac{r^2}{2}m(m+1).
\]
従ってSerre作用を$T^{-r}$とする整数Euler格子としては，unipotence，unimodularity，
rank $3$と任意に大きい$r$が両立する．
\end{proposition}
\begin{proof}
$T$は列$e,Te,T^2e$上の三次差分関係$T^3e=e-3Te+3T^2e$で定義される．
$q(k)=1+k(k+r)/2$は整数値二次多項式で，$q(-k-r)=q(k)$である．
$G_{ij}=q(j-i)$だから差分関係から全ての$\chi(T^ie,T^je)=q(j-i)$が従い，
平行移動不変性とこの対称性が二つの行列等式を与える．行列式は展開して$1$である．
\end{proof}

この模型では$p=(1,-2,1)^{\mathsf t}$はprimitiveなpoint-like類で，$\rho(a,b,c)=a+b+c$である．
$e_1=(0,1,0)^{\mathsf t}$とし，$H$を$e_1-e$のNS像とすると$H^2=1$，$K=-rH$である．
従って幾何的surface-like格子の既存の意味でも
\[
 \delta(G_r)=r^2-9
\]
となる．$r\ge5$の模型は非正defectの仮定を満たさない．
この対応により，模型を実際の非可換Fano曲面へ持ち上げる前に欠けている条件の一つが正確になる．
例えば$r=5$では$q(-1)=-1$であり，負の窓のEuler消滅は失敗する．
この模型はabel圏，ample pair，正の対象，実際のSerre関手を構成したものではない．
従って非可換Kobayashi--Ochiai型定理への反例ではなく，数値的公理だけでは不足することを示す障害である．

\section{三次元のSerre対称性と整数値指数障害}\label{sec:threefold}
ここでは三次元の実際のFano幾何の分類を仮定せず，数値軌道の対称性だけを使う．
以下の数値式を満たす圏の実現は主張しない．

\begin{theorem}[三次元の正規化根とEuler次数]\label{thm:threefold-arithmetic}
$\C$線形Ext-finite abel圏$\mathcal A$の有界導来圏にSerre関手$\SSer$があり，
$\tau=\SSer^{-1}[3]$はheart上exactとする．$\mathcal O$をexceptional対象，
$s$をheart-exactな$r$乗根$s^r=\tau$とし，
$(s_*-I)^M[\mathcal O]=0$をある$M>0$で仮定する．
さらに，$m\gg0$で$\chi(\mathcal O,\tau^m\mathcal O)$が
最高次係数$d/6>0$の有理係数三次多項式$F(m)$に一致するとする．
このとき全整数$m,k$について
\begin{align}
 \chi(\mathcal O,\tau^m\mathcal O)
   &=(2m+1)\left(1+\frac d{12}m(m+1)\right),\label{eq:threefold-anti-euler}\\
 q(k):=\chi(\mathcal O,s^k\mathcal O)
   &=1+\frac{2k}{r}+\frac d{12r^3}k(k+r)(2k+r).\label{eq:threefold-root-euler}
\end{align}
$A=d/r^3$とすると，整数値性の必要十分条件は
\begin{equation}\label{eq:threefold-congruences}
 \boxed{\ A\in\Z,\qquad \frac{A(r+2)}2\in\Z,
 \qquad \frac2r+\frac{A(r+1)(r+2)}{12}\in\Z.\ }
\end{equation}
ここで「十分」は式\eqref{eq:threefold-root-euler}が全整数で整数値になることについてであり，
圏論的な根の存在についてではない．
\end{theorem}
\begin{proof}
Serre双対性と$\SSer=\tau^{-1}[3]$から
\[
 \chi(\mathcal O,\tau^m\mathcal O)
 =-\chi(\tau^m\mathcal O,\tau^{-1}\mathcal O)
 =-\chi(\mathcal O,\tau^{-m-1}\mathcal O)
\]
が全整数$m$で成立する．補題\ref{lem:polynomial}により
$q(k)=F(k/r)$は全整数$k$で成立するので，$F$も全整数$m$で上の反対称性を満たす．
従って$F(m)$は$m\mapsto-m-1$について反対称な三次多項式である．
その一般形は$(2m+1)(a m(m+1)+b)$であり，$F(0)=\chi(\mathcal O,\mathcal O)=1$
より$b=1$，最高次係数から$a=d/12$である．これで二つの式が従う．

任意の三次多項式$q\in\Q[k]$は
$q(0)$と$\Delta q(0)$，$\Delta^2q(0)$，$\Delta^3q(0)$が整数ならば，かつそのときに限り
全整数で整数値となる．実際，Newton展開
$q(k)=\sum_{j=0}^3\Delta^jq(0)\binom{k}{j}$を用いる．
式\eqref{eq:threefold-root-euler}を代入すれば
\[
 \Delta^3q(0)=A,\qquad
 \Delta^2q(0)=\frac{A(r+2)}2,\qquad
 \Delta q(0)=\frac2r+\frac{A(r+1)(r+2)}{12}
\]
となり，所望の条件を得る．
\end{proof}

\begin{corollary}[Euler次数72以下の三次元数値候補]\label{cor:threefold-small-degrees}
前定理で$1\le d\le72$なら，$r>1$について
\[
 \begin{array}{c|c}
 r & d\ \text{の可能な値}\\ \hline
 2 & 8,16,24,32,40,48,56,64,72\\
 3 & 54\\
 4 & 64
 \end{array}
\]
以外の可能性はない．この範囲で$r\ge5$は不可能である．
\end{corollary}
\begin{proof}
$r^3\mid d$だから$r\le4$．$r=2$では式\eqref{eq:threefold-congruences}は
$d/8\in\Z$に帰着する．$r=3$では$A=d/27$は偶数で，
第三条件を合わせると$A\equiv2\pmod6$となり，上界から$A=2$のみである．
$r=4$では$A=d/64$は奇数で，上界から$A=1$のみである．
\end{proof}

\begin{remark}[非可換minifoldとの比較に必要な照合]\label{rem:threefold-minifold}
Loginov--Orlov\cite{LO}の三次元minifoldはEuler次数
$22,40,54,64,72$の五つの数値型を持つ．さらに別論文
\cite{LOM72}で次数72の非可換族が構成されている．
仮に同論文のEuler次数が上の$d$と一致し，かつexceptional構造対象・heart-exactな根・
軌道上の正規化・三次Euler多項式の各仮定が確認できれば，
数値上，次数22では$r>1$はなく，40と72では$r=2$，54では$r=3$，64では$r=2,4$だけが
非自明な候補となる．
これは現時点では\emph{条件付きの照合表}であり，各圏にそれらの根が存在する，
または存在しないと結論したものではない．特に上の$d$と文献中のEuler degreeの定義の一致は
別途確認を要する．
\end{remark}

\section{次数9の認識への適用}
以下でいう六公理は，共同研究原稿\cite[Definition 1.1]{Companion}の条件である．
すなわちnoetherian Ext-finiteな$\C$線形abel圏，$\SSer[-2]$のheart-exact性，
Gabriel filtrationのgeneric quotientでrank oneとなるtorsion-freeな構造対象，その例外性，
$(\mathcal O,\tau)$のAZ ample性，および
$\dim\Hom(\mathcal O,\tau^m\mathcal O)=1+9m(m+1)/2$を課す．
これに$s^3=\tau$を加える．同原稿の基礎補題によりhomological dimensionは2であり，
$A$は右noetherian domain，$\Hom(\mathcal O,s^k\mathcal O)=0$ for $k<0$である．
本節の依存する既存補題は文献\cite{Companion}のラベルで明示する．

\begin{proposition}[最初の障害の消滅]\label{prop:first-obstruction}
上の六公理とexactな三乗根について，軌道上の正規化と反標準軌道上の漸近高次Ext消滅を仮定する．
このとき
\[
 q(k)=\binom{k+2}{2},\qquad E_{-1}=E_{-2}=0.
\]
従って$(\mathcal O,s\mathcal O,s^2\mathcal O)$はexceptional collectionである．
この段階ではstrong性もfull性も結論しない．
\end{proposition}
\begin{proof}
定理\ref{thm:divisibility}に$d=9,r=3$を代入する．
$k=-1,-2$ではHomは零であり，Serre双対性による
$\Ext^2(\mathcal O,s^k\mathcal O)=\Hom(\mathcal O,s^{-k-3}\mathcal O)^*$も零である．
$q(-1)=q(-2)=0$なので中間Extも消える．
\end{proof}

\begin{theorem}[三次normal境界における全消滅]\label{thm:cubic-vanishing}
六公理，exactな三乗根，軌道上の正規化を仮定する．
さらに$0\ne g\in A_3$がnormalであり，$B=A/(g)$がdomainであるとする．
このとき全整数$k$について$E_k=0$であり，
\[
 \dim A_k=\binom{k+2}{2}\ (k\ge0),\qquad
 \dim B_k=3k\ (k>0),\qquad B=\widehat B
\]
が成立する．最後の等号は境界の全section環への自然写像についてである．
三対象はstrong exceptionalとなるが，一般の境界に対してfull性はまだ主張しない．
\end{theorem}
\begin{proof}
まず\cite[\texttt{prop:normal-vanishing}]{Companion}から漸近消滅を得る．
その機構は，AZ2により正の切断の左作用が$E_k$上で高次数では単射になる一方，
$E=H^2_{\mathfrak m}(A)$へのnormal元の左作用は，右torsion作用とnormal twistの自然性により局所冪零となることである．
単射性と局所冪零性から$E_k=0$ for $k\gg0$となる．
従って前命題を適用できる．

$h_k=\dim A_k$，$b_k=\dim B_k$，$P(k)=(k+1)(k+2)/2$とする．
normal切断の完全列から$b_k=h_k-h_{k-3}$，特に$b_{3m}=9m$である．
代数閉体上のdomainの有限次元部分空間$U,V$について
$\dim UV\ge\dim U+\dim V-1$が成立する．
実際，積写像の核の射影化はSegre多様体と交わらず，射影次元定理がこの不等式を与える．
$B_k\ne0$なら三重積に適用して
\[
 3b_k-2\le\dim(B_kB_kB_k)\le b_{3k}=9k,
\]
従って整数性から$b_k\le3k$である．$B_k=0$なら同じ上界は自明である．
$b_0=1$と$h_k=\sum_{j\ge0}b_{k-3j}$から，各剰余類で$h_k\le P(k)$を得る．

一方，$k\ge0$では最上位Extが零なので，Euler多項式から
$h_k=P(k)+\dim E_k\ge P(k)$である．よって$h_k=P(k)$，$E_k=0$ for $k\ge0$．
前命題で$E_{-1}=E_{-2}=0$，さらに$E_k^*=E_{-k-3}$により全整数の消滅を得る．
差分から$b_k=3k$となる．最後にnormal境界のsection完全列
\[
 0\to B_k\to\widehat B_k\to E_{k-3}\to E_k\to0\qquad(k>0)
\]
（\cite[\texttt{cor:cubic-boundary-duality}]{Companion}）により正次数で飽和する．
次数0と負次数も同じnormal切断の長完全列とHomの負次数消滅から従う．
\end{proof}

\begin{corollary}[通常の曲線境界を持つ正規化coindex零の認識]\label{cor:recognition}
前定理の設定で，さらに
\[
 (\qgr B,\pi B,(1))\simeq
 (\coh C,\mathcal O_C,\mathcal L\otimes\alpha^*(-)),
 \quad p_a(C)=1,\quad\deg\mathcal L=3
\]
という通常のintegral curveのtriple表示があれば，$A$は三次元二次Koszul AS-regularである．
既存原稿の高Veronese曲線が滑らかな場合も，その境界復元定理を介して適用できる．
\end{corollary}
\begin{proof}
$B=\widehat B$なので次数3のgenus-one曲線のfull THCRとなり，これは一次生成である．
既存原稿の一次生成三次normal境界の認識定理
\cite[\texttt{thm:cubic-generated-boundary-recognition}]{Companion}を適用する．
高Veroneseからの主張には同原稿\texttt{thm:v10-smooth-veronese-heart}を追加で用いる．
これらの認識・復元定理は本稿の数値論とは独立の入力である．
\end{proof}

\begin{remark}[次数9を所与としない認識への接続]\label{rem:degree-free-route}
次数9のHilbert条件を外したFano曲面候補について，generic rankと整合するpoint-like類を持つ
geometric surface-like数値格子，非正defect，構造対象のrank-one例外類，正規化三乗根が得られれば，
系\ref{cor:defect-index}から数値格子とEuler軌道は射影平面型となる．
従って次の核心はこの\emph{数値曲面性とdefectの符号}をより弱いheartの公理から導くこと，
および得られたEuler窓を実際の消滅・full性へ持ち上げることである．
この方向は次数上界を単に定義に加えるより構造的だが，現時点では追加条件付きの経路である．
\end{remark}

\section{現在の結論と未解決の範囲}
滑らかなdel Pezzo曲面では非正規化exact指数の次数別最大値が反標準次数$1,2,3,4,5,6$を達成する一方，次数7では1，次数8では4，次数9では3になる（定理\ref{thm:delpezzo-exact-index-table}）．
一般の次数4 del Pezzo曲面はexact指数2だが，同じ数値格子には四乗根が存在する．これは数値根がexact自己同値へ持ち上がらない具体例である（命題\ref{prop:degree-four-lifting-gap}）．
この分類は古典的自己同型を組み合わせたもので，非可換曲面の実現やfullnessを結論しない．
軌道が数値群を張る場合には全体と軌道の正規化が一致する（命題\ref{prop:orbit-spanning-normalization}）．

正規化は古典的Fano指数を忠実に回復する．全数値群と一つの軌道の正規化は同じ次数集合を与えるが，許される根と標識依存性は異なる．
surface-likeな擬格子の枠組みで，数値根を標準類のdivisibilityとして識別した．
非正defectを加えた数値分類は既知のKuznetsov分類の帰結であり，新しい曲面分類とは主張しない．
曲面ではEuler多項式の整数値性から
$r^2\mid d$と偶数次数の追加因子2が従い，$1\le d\le9$のcoindex零は次数9へ帰着する．
しかし$d\le9$自体，漸近消滅，負の窓の消滅，およびfull性は一般公理から自動ではない．
数値模型はこれらを区別する必要を示す．
三次normal境界では正規化により有限section欠損が消えるが，特異な高Veronese境界の
原始剰余層の可逆性と，元の境界の一次生成性は別に残る．
normal切断がない場合には，漸近消滅と生成障害の零性を優先して調べる．

本稿の証明は主として数値線形代数と既知の再構成結果の組合せである．
独立した新規性は未確定であり，一般の非可換射影空間認識を完了したとは主張しない．
付属\texttt{verify\_normalized\_index.py}は整数値条件，模型のSerre行列，正規化根の行列式，
反例の数値作用，数値的tensorの群法則，標準類の整除性，defect模型を厳密算術で検算する．それは抽象圏の存在や一般証明の形式的認証ではない．

\begingroup\small
\begin{thebibliography}{99}
\bibitem{KS} A.~Kuznetsov and M.~Smirnov,
\emph{On residual categories for Grassmannians}, arXiv:1802.08097v3 (2019).
\url{https://arxiv.org/abs/1802.08097}. Definition 2.1, Example 2.2.
\bibitem{BO} A.~Bondal and D.~Orlov,
\emph{Reconstruction of a variety from the derived category and groups of autoequivalences},
arXiv:alg-geom/9712029v1 (1997).
\url{https://arxiv.org/abs/alg-geom/9712029}. Theorem 3.1.
\bibitem{KK} Y.~Kachi and J.~Koll\'ar,
\emph{Characterizations of $\mathbb P^n$ in arbitrary characteristic},
arXiv:math/0003230v1 (2000).
\url{https://arxiv.org/abs/math/0003230}. Corollary 2(1).
\bibitem{Keeler} D.~S.~Keeler,
\emph{Criteria for $\sigma$-ampleness}, J. Amer. Math. Soc. \textbf{13} (2000), 517--532.
\url{https://arxiv.org/abs/math/9912051}. arXiv v2, Theorems 1.2--1.3.
\bibitem{dTVdB} L.~de Thanhoffer de Volcsey and M.~Van den Bergh,
\emph{On an analogue of the Markov equation for exceptional collections of length 4},
arXiv:1607.04246v1 (2016). \url{https://arxiv.org/abs/1607.04246}.
\S3.2--3.3; surface-type Serre latticeとnumerical codimension filtration.
\bibitem{Kuz} A.~Kuznetsov,
\emph{Exceptional collections in surface-like categories},
Sb. Math. \textbf{208} (2017), 1368--1398.
\url{https://arxiv.org/abs/1703.07812}.
\url{https://doi.org/10.1070/SM8917}.
\S3, Definition 4.3, Definition 5.3, Theorem 5.12（出版版）．
arXiv v2では最後の定理は5.11であり，defect節の一部の番号も異なる．
\bibitem{MU} I.~Mori and K.~Ueyama,
\emph{A categorical characterization of quantum projective spaces},
arXiv:1708.00167v4 (2021). \url{https://arxiv.org/abs/1708.00167}.
Theorem 4.1; v1の誤った証明に依存しない．
\bibitem{Hanihara} N.~Hanihara,
\emph{Calabi--Yau completions for roots of dualizing dg bimodules},
arXiv:2412.18753v1 (2024). \url{https://arxiv.org/abs/2412.18753}.
Definition 2.5, Definitions 3.1--3.2, Theorem 3.3.
\bibitem{LO} K.~Loginov and D.~Orlov,
\emph{Non-commutative Fano varieties and three-dimensional minifolds},
arXiv:2609.39766 (2026). \url{https://arxiv.org/abs/2609.39766}.
Definitions 2.1, 2.3, 2.6; Proposition 2.7（2026年11月17日に取得した一次本文）．
\bibitem{LOM72} K.~Loginov and D.~Orlov,
\emph{Three-dimensional noncommutative Fano minifolds},
arXiv:2610.31250v1 (2026). \url{https://arxiv.org/abs/2610.31250}.
\bibitem{DolgachevMartin} I.~Dolgachev and G.~Martin,
\emph{Automorphisms of del Pezzo surfaces in odd characteristic},
J. London Math. Soc. \textbf{109} (2024), e12905.
\url{https://doi.org/10.1112/jlms.12905}.
古典的del Pezzo分類と自己同型群（奇標数の記述を複素数体の例の背景にも使用）．
\bibitem{BauerCatanese} I.~Bauer and F.~Catanese,
\emph{$\mathfrak S_5$-equivariant syzygies for the del Pezzo surface of degree 5},
Rend. Circ. Mat. Palermo (2) \textbf{70} (2021), 97--120.
\url{https://doi.org/10.1007/s12215-020-00483-9}.
\S2; degree $5$の自己同型$\mathfrak S_5$とpentagonの幾何．
\bibitem{Rieken} S.~Rieken,
\emph{Moduli of real pointed quartic curves},
Geom. Dedicata \textbf{185} (2016), 171--203.
\url{https://doi.org/10.1007/s10711-016-0174-0}.
\S3.3; degree $2$のGeiser involutionとPicard作用の背景．
\bibitem{HassettKreschTschinkel} B.~Hassett, A.~Kresch, and Y.~Tschinkel,
\emph{On the moduli of degree $4$ del Pezzo surfaces},
Development of Moduli Theory---Kyoto 2013,
Adv. Stud. Pure Math. \textbf{69} (2016), 349--386.
\url{https://doi.org/10.2969/aspm/06910349}.
Introduction, \S2; 一般の次数4曲面の自己同型群$(\Z/2)^4$とrank $4$ quadric pencils.
\bibitem{Milne} J.~S.~Milne,
\emph{Class Field Theory}, version 4.03 (2020).
\url{https://www.jmilne.org/math/CourseNotes/CFT.pdf}.
Chapter II, Example 1.20（巡回群の$H^1$とnorm核）．
\bibitem{Companion} 本研究の現行原稿，
\emph{反Serre作用の三乗根による非可換射影平面の認識},
\texttt{nc\_dp9\_recognition.tex}, 2026年11月17日，現行原稿（認識定理の数学的内容は維持）．
参照は現行ソースの定理ラベルで固定する．本稿で使用する仮定と結論を各箇所で明示した．
\end{thebibliography}
\endgroup
\end{document}
